\section{Conclusion and Future Directions}
\label{sec:conclusion}

In this paper, we introduced \textbf{Unitary Geodesic Routing (UGR)}, a novel test-time model ensembling framework that rejects the long-standing assumption that routing states must be updated in flat Euclidean spaces. By modeling the ensembling state on the curved $(K-1)$-dimensional hypersphere $\mathbb{S}^{K-1}$ and mapping coordinates to the probability simplex via Born's rule, UGR achieves a Softmax-free projection that guarantees native simplex constraints with zero geometric distortion or activation scale-mismatch. We derived an efficient, closed-form Rodrigues-like geodesic rotation operator that interpolates along the shortest great-circle path of the hypersphere without needing numerical solvers or matrix exponentials. To solve the stability-plasticity dilemma, we introduced Torque-Driven Adaptive Agility, which dynamically scales the update step size based on representational torque (angular mismatch), and Spatial-Temporal Geodesic Coupling to smoothly propagate ensembling trajectories across query boundaries. Evaluated across 10 random seeds in the high-fidelity 14-layer Analytical Coordinate Sandbox (ICS), UGR achieves a Joint Mean Accuracy of \textbf{75.08\%} (outperforming ChemMerge Reset by \textbf{+5.43\%}) and slashes layer-to-layer routing jitter ($L \ge 5$) by \textbf{2.10$\times$}.

\textbf{Future Directions.} There are several exciting paths for future research:
\begin{enumerate}
    \item \textbf{Differentiable Training-Time Geodesic Routing:} In this work, UGR is applied training-free at test-time. A natural next step is to formulate geodesic flow as a differentiable layer for end-to-end training. Since all operations in UGR—including the Born's rule mapping ($\alpha_k = s_k^2$), the orthonormal basis construction ($\mathbf{u}$), and the geodesic rotation ($\mathbf{s}_t^{(l)} = \cos(\theta)\mathbf{s}_t^{(l-1)} + \sin(\theta)\mathbf{u}$)—are formulated using standard differentiable trigonometric, algebraic, and normalization operators, gradients can be seamlessly backpropagated through our geometric router. Specifically, the gradient of any loss function $\mathcal{L}$ with respect to the previous state $\mathbf{s}_t^{(l-1)}$ and the target $\mathbf{w}_t^{(l)}$ can be derived analytically using standard automatic differentiation frameworks. This allows us to optimize the adapter task centroids $\{\mu_k^{(l)}\}$ or learn the step size $\eta$ jointly with the model backbone under a unified geometric optimization objective.
    \item \textbf{Real-World Validation and Scaling to Large Language Models:} While our 14-layer Analytical Coordinate Sandbox (ICS) provides a highly controlled, high-fidelity environment to mathematically isolate and audit routing dynamics across 10 independent seeds, we have successfully validated UGR's real-world efficacy on a multi-task text classification Mixture-of-Experts benchmark. Our results on the \texttt{20newsgroups} block-structured stream show that UGR translates seamlessly to real text representations, achieving a massive \textbf{+21.60\%} absolute accuracy boost while reducing routing jitter by \textbf{16.8$\times$} compared to ChemMerge. An exciting next step is to scale UGR to ensemble token-level LoRA expert adapters in Large Language Models (such as LLaMA-3 and Mistral) across non-stationary, sequential multi-task text generation streams (e.g., transitioning from mathematical reasoning to creative writing, or processing alternating prompts from GLUE benchmarks, MMLU subcategories, and code-generation datasets). Explicitly, evaluating UGR on real token-by-token autoregressive decoding sequences represents our immediate next-step experimental milestone to align the empirical scope with our LLM serving blueprint. Additionally, we are preparing to evaluate UGR on Vision Transformers (ViT) with task-specific adapters under streaming image classification. The analytical Jacobians and positive orthant persistence proofs derived in Appendix \ref{appendix:formulations} provide the complete mathematical foundation for this end-to-end differentiable, real-world integration, positioning UGR as a highly scalable alternative for production-grade test-time adaptive serving.
    \item \textbf{Alternative Lie Groups and Symmetric Spaces:} While this work focuses on the unit hypersphere $\mathbb{S}^{K-1}$, we aim to extend this geometric paradigm to more general Lie groups and symmetric spaces, exploring how different non-Euclidean manifold geometries influence multi-task representations. A highly promising and forward-looking direction is to utilize the Stiefel manifold $\text{St}(p, D)$—the space of $p$-dimensional orthonormal frames in $\mathbb{R}^D$—or the Grassmannian manifold $\text{Gr}(p, D)$ to directly ensemble parameter matrices (such as the low-rank projection weights of LoRA adapters) rather than just routing weights. By performing geodesic flows directly on these matrix manifolds, we could continuously interpolate adapter weights along the shortest curved geodesic paths of the symmetric space, guaranteeing that the blended parameters maintain exact orthonormality and structural ranks, thereby bypassing the heuristic weighted average blending that typically degrades model capacity.
    \item \textbf{Online Centroid Estimation and Adaptation under Semantic Drift:} In this work, the expert task centroids $\boldsymbol{\mu}_k^{(l)}$ are pre-computed during a calibration phase and assumed to be static. However, real-world serving streams often experience gradual semantic drift or introduce out-of-distribution domains. An extremely valuable future direction is to explore test-time online clustering or centroid adaptation. Specifically, we can formulate an online, exponentially decaying centroid update rule:
    \begin{equation}
    \boldsymbol{\mu}_{k, t}^{(l)} = (1 - \gamma_{k,t})\boldsymbol{\mu}_{k, t-1}^{(l)} + \gamma_{k,t} h_t^{(l-1)}
    \end{equation}
    where $\gamma_{k,t} = \gamma_0 \cdot \alpha_{k, t}^{(l)}$ is a learning rate dynamically scaled by the routing coefficient $\alpha_{k, t}^{(l)}$, and $\gamma_0 \in (0, 1)$ is a base decay rate. This ensures that expert centroids are updated on-the-fly proportionally to their contribution to the current prediction. By dynamically updating centroids as new inputs arrive, UGR can remain robust to long-term semantic drift and adaptively route queries under evolving task boundaries without requiring offline recalibration.
    \item \textbf{Layer-Dependent Step Size Adaptivity:} In our current formulation, the step size scaling parameter $\eta$ is held constant across all network layers. In deep neural architectures, representations at earlier layers are often highly plastic, serving to extract low-level features, whereas deeper layers contain highly structured, specialized semantic features. A promising path is to explore layer-dependent step sizes, where earlier layers utilize a larger $\eta$ (maximizing adaptability and prompt-boundary agility) and deeper layers utilize a smaller $\eta$ (maximizing representation stability and temporal smoothing), which could further optimize the accuracy-stability Pareto frontier.
\end{enumerate}
